\section{Operations Involving the Natural Numbers}
\begin{questions}
    % Set Difference
    \item The set difference operation on a set $A$ can be defined as the function $f: \Power{A}^2 \to \Power{A}$ where
    \[
        f(S, T) = \{x \in A \vert x \in S \land x \notin T\}
    \]
    for all $S, T \in \Power{A}$.

    % Successor of $n$
    \item \begin{proof}
        Note $\Successor{n} = \Successor{n + 0}$ by \textbf{A1}, which is the same as $n + \Successor{0}$ by \textbf{A2}. But since $\Successor{0} = 1$, we see that $\Successor{n} = n + 1$ as needed.
    \end{proof}

    % Sum of $\Successor{d}$ and $n$
    \item \begin{proof}
        We induct on $n$.

        When $n = 0$, the left-hand side is $\Successor{d} + 0 = \Successor{d}$ by \textbf{A1}. The right-hand side similarly simplifies to $\Successor{d + 0} = \Successor{d}$ by \textbf{A1} too. Thus the base case is proven.

        Suppose $\Successor{d} + k = \Successor{d + k}$ for some $k \in \N$. We are to prove that $\Successor{d} + \Successor{k} = \Successor{d + \Successor{k}}$.

        One sees that
        \begin{align*}
            \Successor{d} + \Successor{k} &= \Successor{\Successor{d} + k} & (\text{by \textbf{A2}})\\
            &= \Successor{\Successor{d + k}} & (\text{induction hypothesis})\\
            &= \Successor{d + \Successor{k}} & (\text{by \textbf{A2}})
        \end{align*}
        which proves the case for $\Successor{k}$.
    \end{proof}

    % A Commutative Result
    \item \begin{proof}
        We see that
        \begin{align*}
            (a + b) + c &= a + (b + c) & (\text{by \myref{theorem:addition-associative-naturals}})\\
            &= a + (c + b) &  (\text{by \myref{theorem:addition-commutative-naturals}})\\
            &= (a + c) + b & (\text{by \myref{theorem:addition-associative-naturals}})
        \end{align*}
        which proves the corollary.
    \end{proof}

    % Addition that Equals Self
    \item \begin{proof}
        We only prove the case that if $a + b = b$ then $a = 0$, since the other is easily implied by commutativity of addition (\myref{theorem:addition-commutative-naturals}).
    
        We can write $a + b = b$ as $a + b = b + 0$ by \textbf{A1}. Using commutativity of addition on the right-hand side we have
        \[
            a + b = 0 + b.
        \]
        Thus, using \myref{prop:addition-cancellation-property} we get $a = 0$ as required.
    \end{proof}

    % Proof of Commutativity
    \item \begin{partquestions}{\roman*}
        \item \begin{proof}
            We consider induction on $n$.

            The base case of $n = 0$ is easily proven since the left-hand side, being $0 \times 0$, simplifies to 0 using \textbf{M1}.

            Suppose that $0 \times k = 0$ for some $k \in \N$. We show that $0 \times \Successor{k} = 0$.

            Note that
            \begin{align*}
                0 \times \Successor{k} &= (0 \times k) + 0 & (\text{by \textbf{M2}})\\
                &= 0 \times k & (\text{by \myref{lemma:sum-of-0-and-n}})\\
                &= 0 & (\text{induction hypothesis})
            \end{align*}
            which proves the case for $\Successor{k}$.
        \end{proof}

        \item \begin{proof}
            We induct on $n$.

            When $n = 0$, the left-hand side is $m \times 0$ which is 0 by \textbf{M1}, and the right-hand side is $0 \times m$ which is 0 by \myref{lemma:product-of-0-and-n}. So the base case is proven.

            Suppose $m \times k = k \times m$ for some $k \in \N$. We show $m \times \Successor{k} = \Successor{k} \times m$.

            Note
            \begin{align*}
                m \times \Successor{k} &= (m \times k) + m & (\text{by \text{M2}})\\
                &= (k \times m) + m & (\text{induction hypothesis})\\
                &= \Successor{k} \times m & (\text{by \myref{lemma:product-of-successor-and-n}})
            \end{align*}
            which proves the case for $\Successor{k}$.
        \end{proof}
    \end{partquestions}

    % Product that Equals Self
    \item \begin{proof}
        We only prove the case that if $b \neq 0$ and $a \times b = b$ then $a = 1$, since the other is easily implied by commutativity of multiplication (\myref{theorem:multiplication-commutative-naturals}).
    
        Note $a \neq 0$ since, otherwise, the left-hand side will be 0 (by \myref{lemma:product-of-0-and-n}) and the right-hand side is $b \neq 0$, a contradiction. So we can write $a \times b = b$ as $a \times b = 1 \times b$ by \myref{corollary:multiplication-by-1}. Thus, using \myref{prop:multiplication-cancellation-property} we get $a = 1$ as required.
    \end{proof}

    % Exponentiation is Unlike the Rest
    \item \begin{partquestions}{\alph*}
        \item False. One sees
        \begin{align*}
            0 \wedge 1 &= 0^1\\
            &= (0^0) \times 0 & (\text{by \textbf{E2}})\\
            &= 0 & (\text{by \textbf{M1}})
        \end{align*}
        but $1 \wedge 0 = 1^0 = 1$ by \textbf{E1}.
        \item False. One sees $(2\wedge 3)\wedge2 = \left(2^3\right)^2 = 64$ but $2\wedge (3\wedge 2) = 2^{\left(3^2\right)} = 512$.
    \end{partquestions}

    % Power of One
    \item \begin{proof}
        We see that
        \begin{align*}
            n^1 &= n^\Successor{0}\\
            &= n^0 \times n & (\text{by \textbf{E2}})\\
            &= 1 \times n & (\text{by \textbf{E1}})\\
            &= n, & (\text{by \myref{corollary:multiplication-by-1}})
        \end{align*}
        proving this proposition.
    \end{proof}

    % Product to a Power
    \item \begin{proof}
        We induct on $n$.

        When $n = 0$ we have
        \begin{align*}
            a^0 \times b^0 &= 1 \times 1 & (\text{by \textbf{E1}})\\
            &= 1 & (\text{by \myref{corollary:multiplication-by-1}})\\
            &= (a\times b)^0 & (\text{by \textbf{E1}})
        \end{align*}
        so the base case is proven.

        Suppose $a^k \times b^k = (a\times b)^k$ for some $k \in \N$. We show that $a^\Successor{k} \times b^\Successor{k} = (a\times b)^\Successor{k}$.

        Note that
        \begin{align*}
            a^\Successor{k} \times b^\Successor{k} &= \left(a^k \times a\right) \times \left(b^k \times b\right) & (\text{by \textbf{E2}})\\
            &= \left(a^k \times b^k\right) \times (a \times b) & (\text{properties of multiplication})\\
            &= (a\times b)^k \times (a \times b) & (\text{induction hypothesis})\\
            &= (a\times b)^\Successor{k} & (\text{by \textbf{E2}})
        \end{align*}
        which proves the case for $\Successor{k}$.
    \end{proof}
\end{questions}
